% ============================================================
% 슈퍼마리오 메이커 2 — 전체 파츠 카탈로그
% 프로젝트: 협업형 UGC 플랫포머 (madcamp 26s-w2-c3-08)
% 컴파일: Overleaf에서 Compiler를 XeLaTeX 로 설정할 것
%   (Menu > Settings > Compiler > XeLaTeX)
%
% [작성 근거] 2026-07-10 사용자 제공 위키 원문 3종:
%   1차: Super Mario Wiki (mariowiki.com) "Super Mario Maker 2" Course parts 절
%   2차: Super Mario Maker 2 Wiki (Fandom) Course Parts / Terrain / Items /
%        Enemies / Gizmos 페이지 (파츠 총수·카테고리 검증용)
%   3차: Kaizo Mario Maker Wiki "SMM2: List of Course Elements"
%        (파츠 개수 제한, 에디터 설정)
%   카테고리 소속은 mariowiki(v3.0.0 기준)를 따른다.
%   단, "필요 애니메이션" 열은 원문이 아니라 팀의 자체 분석이다.
% ============================================================
\documentclass[10pt,a4paper]{article}
\usepackage{kotex}                % 한국어 (XeLaTeX)
\usepackage[margin=2cm]{geometry}
\usepackage{longtable}
\usepackage{booktabs}
\usepackage{array}
\usepackage{xcolor}
\usepackage[colorlinks=true,linkcolor=blue!60!black,urlcolor=blue!60!black]{hyperref}
\usepackage{enumitem}
\setlist{nosep}

\newcolumntype{L}[1]{>{\raggedright\arraybackslash}p{#1}}

\title{슈퍼마리오 메이커 2 전체 파츠 카탈로그\\
\large 공식 카테고리 기준 분류 및 파츠별 특성·애니메이션 명세 (v3.0.0 기준)}
\author{madcamp 2주차 팀 26s-w2-c3-08}
\date{2026년 7월 10일}

\begin{document}
\maketitle

\begin{quote}\small
\textbf{작성 근거:} 파츠 목록(2--5절)과 특성 서술은 Super Mario Wiki
(mariowiki.com)의 Course parts 절 원문을 근거로 하며, 파츠 총수와 카테고리
구성은 SMM2 Fandom 위키(지형 17종 / 아이템 18종 / 적 44종 / 장치 40종,
합계 119종)로 교차 검증했다. ``필요 애니메이션'' 열은 원문이 아닌
\textbf{팀 자체 분석}이다. 스타일 표기: \textbf{전체}=5스타일,
\textbf{클래식}=SMB1/SMB3/SMW/NSMBU, \textbf{3DW}=3D월드 전용.
버전 주석: [2.0]=v2.0.0 추가, [3.0]=v3.0.0 추가.
\end{quote}

\tableofcontents
\newpage

% ============================================================
\section{파츠 기술 기준 (스키마)}
% ============================================================

모든 파츠는 아래 기준(필드)으로 기술한다. 이 스키마는 그대로 우리 프로젝트의
에셋 데이터 모델 초안으로 사용할 수 있다.

\begin{longtable}{L{3cm} L{12.5cm}}
\toprule
\textbf{필드} & \textbf{설명} \\
\midrule
이름 & 한국어 명칭 / 영문 명칭 \\
스타일 & 사용 가능한 게임 스타일 (전체 / 클래식 / 3DW / 개별) \\
특성 & 자체 동작·플레이어 상호작용 (위키 원문 근거) \\
변형·비고 & 대체 형태(alternate form), 색상 변형, 해금 조건 등 \\
애니메이션 & 표현에 필요한 애니메이션 상태 (\ref{sec:anim}절 어휘, 팀 분석) \\
\bottomrule
\end{longtable}

% ------------------------------------------------------------
\subsection{수식자(modifier) 체계}\label{sec:mod}
% ------------------------------------------------------------

원문의 Modifier legend에 따르면, 파츠에는 아래 12종의 수식자·배치 옵션이
파츠별로 선택적으로 적용된다. \textbf{파츠 종류를 늘리지 않고 조합으로
다양성을 만드는 장치}로, 우리 속성 시스템의 직접적 선례다.

\begin{enumerate}
  \item 슈퍼버섯 부착 (거대화 --- 굼바는 밟으면 2마리로 분열, 엉금 등껍질은 블록 파괴 등 개체별 효과)
  \item 날개 부착 (Wings --- 개체별로 비행·도약 등 행동 변화)
  \item 낙하산 부착 (Parachute)
  \item 열쇠 부착 (처치 시 열쇠 드롭)
  \item 다른 파츠와 적층(stack) 가능 (3DW 제외)
  \item 토관·킬러 대포에 넣어 배출 가능
  \item 블록(벽돌/?/숨김/음표 등)에 수납 가능
  \item 라키투가 투척하는 탄으로 지정 가능 (3DW 제외)
  \item 라키투 구름·클라운 카에 탑승시켜 조종수로 지정 가능 (3DW 제외)
  \item 나무(Tree)에 수납 가능 (3DW 전용)
  \item 중간 깃발에 파워업 수납 가능
  \item 레일(Track)에 부착해 이동 가능 (3DW 제외)
\end{enumerate}

파츠별 수식자 적용 가능 여부 전수 매트릭스는 분량 관계로 본 문서에서
생략했다 (원문에 존재하므로 필요 시 추가 가능).

% ------------------------------------------------------------
\subsection{애니메이션 어휘 정의}\label{sec:anim}
% ------------------------------------------------------------

우리 프로젝트에서 유저 생성 애니메이션은 \textbf{idle / walk / onair} 3종으로
확정되었으므로, 그 외 어휘는 코드 연출(절차적)로 구현하거나 생략 대상이다.

\begin{longtable}{L{2.5cm} L{7cm} L{6cm}}
\toprule
\textbf{어휘} & \textbf{의미} & \textbf{우리 프로젝트 구현} \\
\midrule
static & 애니메이션 없음 (정지 이미지 1장) & 이미지 1장 \\
idle & 제자리 대기 루프 & AI 생성 (확정) \\
walk & 지상 이동 루프 & AI 생성 (확정) \\
onair & 공중 상태 루프 (점프·낙하·비행 공용) & AI 생성 (확정) \\
swim & 수영 루프 & 물 스테이지 도입 시에만 \\
attack & 공격 동작 (투척·발사) & 미도입 권장 (속성으로 대체) \\
death / squash & 사망·밟힘 연출 & 코드 연출 \\
bump / shake & 블록 튕김·낙하 예고 떨림 & 코드 연출 (트윈) \\
spin / rotate & 회전 루프 & 코드 연출 (이미지 회전) \\
toggle & ON/OFF 등 상태 전환 & 코드 연출 (이미지 교체) \\
emerge & 출현 (토관·땅에서 등장) & 코드 연출 (마스크+이동) \\
경로 이동 & 레일·경로 추종 & 코드 연출 (트윈) \\
\bottomrule
\end{longtable}

\newpage
% ============================================================
\section{지형 (Terrain)}
% ============================================================

벽돌 블록 아이콘, 청록색 카테고리. 총 17종(별도 파츠 기준).
전부 static이 기본이며 상호작용 연출(bump 등)은 코드 처리 대상이다.

\begin{longtable}{L{2.5cm} L{1.3cm} L{6.2cm} L{3.2cm} L{1.9cm}}
\toprule
\textbf{파츠} & \textbf{스타일} & \textbf{특성} & \textbf{변형·비고} & \textbf{애니메이션} \\
\midrule
\endhead
땅 Ground & 전체 & 파괴 불가 기본 지형. 플레이어·적이 보행, 파츠 배치 가능 & — & static \\
급경사 Steep Slope & 전체 & 45°(1:1) 경사. 미끄럼 내려가기(적 처치) 가능. 방향·길이 조절 & 계단형 변형(기능 동일) & static \\
완경사 Gentle Slope & 전체 & 2:1 비율 완경사. 기능은 급경사와 동일 & 계단형 변형 & static \\
토관 Pipe & 전체 & 서브 에어리어 워프 출입구 또는 파츠 배출구. 색상별 배출 속도: 파랑$<$초록$<$노랑$<$빨강 & 색 4종 & static (배출은 emerge) \\
투명 토관 Clear Pipe & 3DW & 내부 이동이 보이는 관. 연장·굴절 가능. 적·아이템도 통과 & — & static \\
가시 Spike Trap & 클래식 & 접촉 피해. 구두·요시·드라이본 등껍질 착용 시 안전 보행 & 수중에서 젤렉트로(SMB3)·성게(SMW)로 외형 변경 & static \\
버섯 발판 Mushroom Platform & 클래식 & 크기 조절 가능한 반통과 발판 & 빨강/초록/노랑 & static \\
반통과 발판 Semisolid Platform & 전체 & 아래에서 통과, 위에 착지. 3DW에서 고양이 마리오가 벽면 오르기 가능 & 외형 3종 & static \\
다리 Bridge & 클래식 & 가로 길이 조절 반통과 발판 & — & static \\
블록 Block & 전체 & 강한 형태로 아래서 치면 파괴. 아이템·적 수납 가능. 코인 수납 시 연타로 최대 10코인 & SMW: 회전 블록(타격 시 회전, 일시 비충돌). 3DW 지하·숲: 크리스탈 블록 & static + bump \\
?블록 ? Block & 전체 & 기본 코인 1개 수납, 타격 후 빈 블록화. 다른 파츠 수납 가능 & 빈 블록(Empty Block) & idle(반짝임) + bump \\
딱딱한 블록 Hard Block & 전체 & 통상 파괴 불가. 폭탄병 폭발·빅 마리오 점프로 파괴 가능 & 3DW: 록 블록 & static \\
숨겨진 블록 Hidden Block & 전체 & 비가시. 아래에서 치면 실체화. 수납 가능 & — & static \\
도넛 블록 Donut Block & 전체 & 일정 시간 서 있으면 낙하. 쿵쿵 착지로도 작동. 3DW는 완전 충돌, 클래식은 반통과 & — & static + shake + 낙하 \\
음표 블록 Note Block & 클래식 & 밟거나 치면 크게 튕김. 수납 가능 & — & bump \\
음악 블록 Music Block & 클래식 & 음표 블록의 변형. 밟은 대상별 악기, 배치 높이별 음정 재생 & — & bump \\
구름 블록 Cloud Block & 전체 & 반통과 발판. 쿵쿵·스큐어·쿠파(SMB3)·모튼 GP로 파괴됨. SMB2 마리오는 파고들기 가능 & — & static \\
얼음 블록 Ice Block & 전체 & 저마찰(미끄러움). 강한 힘으로 파괴 가능 & — & static \\
\bottomrule
\end{longtable}

% ============================================================
\section{아이템 (Items)}
% ============================================================

코인 아이콘, 마젠타색 카테고리. 총 18종(별도)·32종(고유 명칭).
파워업의 본질은 플레이어 상태 변경이므로 효과 로직이 핵심이다.

\begin{longtable}{L{2.9cm} L{1.5cm} L{6.6cm} L{2.6cm} L{1.5cm}}
\toprule
\textbf{파츠} & \textbf{스타일} & \textbf{특성} & \textbf{변형·비고} & \textbf{애니메이션} \\
\midrule
\endhead
코인 Coin & 전체 & 100개당 1UP. 토관·대포·라키투 부착 시 다수 생성 & — & idle \\
얼린 코인 Frozen Coin [2.0] & 클래식 & 얼음 속 코인. 파이어볼·용암방울·파이어바·성난 태양 등 화염으로 융해. P스위치 시 얼린 벽돌로 변함 & 코인의 변형 & static \\
10/30/50코인 & 전체 & 각 10/30/50코인 가치의 대형 코인 & 10코인의 변형 & idle \\
핑크 코인 Pink Coin & 전체 & 구역 내 전부 수집 시 열쇠 출현 & — & idle \\
슈퍼버섯 Super Mushroom & 전체 \textcolor{red}{(3DW 여부 원문 모호)} & 슈퍼 형태로 변신. 적에게 부착 시 거대화 & — & walk(굴러감) \\
마스터 소드 Master Sword [2.0] & SMB1 & 슈퍼버섯의 변형. 링크로 변신: 검·방패·활·폭탄·아래찌르기 & 링크 전용 모션 다수 & idle \\
파이어플라워 Fire Flower & 전체 & 파이어 형태, 파이어볼 발사 & — & idle \\
슈퍼볼 플라워 Superball Flower & SMB1 & 파이어플라워의 변형. 벽 반사 슈퍼볼. 스토리 모드 클리어로 해금 & 전용 BGM(마리오랜드) & idle \\
빅 버섯 Big Mushroom & SMB1 & 거대화. 딱딱한 블록 등 파괴. 피격 시 슈퍼 형태로 강등 & — & walk \\
슈퍼 리프 Super Leaf & SMB3 & 너구리 형태 (비행·꼬리치기) & — & idle \\
망토 깃털 Cape Feather & SMW & 망토 형태 (활공·망토치기) & — & idle \\
프로펠러 버섯 Propeller Mushroom & NSMBU & 프로펠러 형태 (급상승·활강) & — & idle \\
슈퍼 도토리 Super Acorn [3.0] & NSMBU & 날다람쥐 형태 (활공·벽 붙기) & — & idle \\
파워 벌룬 Power Balloon [3.0] & SMW & 풍선 형태로 공중 부유. 피격 전까지 유지 & — & idle \\
SMB2 버섯 SMB2 Mushroom [3.0] & SMB1 & SMB2 물리로 변경: 적 위에 타기·들기·던지기, 차지 점프 & — & walk \\
개구리 슈트 Frog Suit [3.0] & SMB3 & 수영 강화, 수면 위 달리기(P미터 조건) & — & idle \\
슈퍼 벨 Super Bell & 3DW & 고양이 형태: 벽·반통과 발판·골 기둥 오르기, 할퀴기 & — & idle \\
슈퍼 해머 Super Hammer & 3DW & 빌더 형태: 해머 공격, 빌더 박스 설치. 스토리 모드 클리어로 해금 & — & idle \\
부메랑 플라워 Boomerang Flower [3.0] & 3DW & 부메랑 형태 & — & idle \\
슈퍼스타 Super Star & 전체 & 일시 무적 & — & walk + 점멸 \\
1UP 버섯 & 전체 & 잔기 +1 & 밤 지상에서 썩은 버섯으로 대체 & walk \\
썩은 버섯 Rotten Mushroom & 클래식 (밤 지상) & 플레이어를 추적, 접촉 시 피해 & 1UP 버섯의 변형 & walk \\
구두 굼바 Shoe Goomba / 하이힐 Stiletto & SMB1·SMB3 & 구두 탄 굼바. 처치 후 구두 탑승: 가시·검은 뻐끔 위 안전. 빅 슈는 착지 스톰프, 하이힐은 쿵쿵·킬러 대포도 밟기 가능, 날개 시 플러터 점프 & 하이힐은 변형 & walk \\
요시 알 Yoshi's Egg & SMW·NSMBU & 지면 접촉 시 요시 부화. 탑승·삼키기(뱉으면 발사체) & 열쇠 부화 설정 가능 & idle $\to$ 부화 \\
빅 빨간 요시 알 & SMW·NSMBU & 빨간 요시 부화: 혀 대신 화염 브레스, 탑승자가 파이어 형태면 3연발 & 요시 알의 변형 & idle $\to$ 부화 \\
착용 상자 5종 [3.0] (대포/프로펠러/굼바 가면/킬러 가면/빨간 POW) & 3DW & 머리 착용: 포격(차지 가능) / 공중 3회 비행 / 적이 공격 안 함 / 수평 비행 / 충격파 3회 & — & static \\
\bottomrule
\end{longtable}

% ============================================================
\section{적 (Enemies)}
% ============================================================

굼바 아이콘, 초록색 카테고리. 총 44종(별도)·52종(고유 명칭)으로 최다.
클라운 카·라키투 구름 등 탑승물도 이 카테고리 소속이다(mariowiki 기준).

\begin{longtable}{L{2.7cm} L{1.4cm} L{6.2cm} L{2.8cm} L{1.9cm}}
\toprule
\textbf{파츠} & \textbf{스타일} & \textbf{특성} & \textbf{변형·비고} & \textbf{애니메이션} \\
\midrule
\endhead
굼바 Goomba (SMW: 굼바돌이 Galoomba) & 전체 & 보행, 절벽에서 낙하. 밟으면 즉사(굼바돌이는 뒤집혀 들기 가능). 3DW에서는 발견 시 추적. 밤 지상에서 공중 부유. 거대화 시 밟으면 2마리로 분열 & Goombrat/Goombud: 절벽에서 방향 전환 & walk, squash \\
엉금엉금 Koopa Troopa & 전체 & 밟으면 등껍질화(SMW·3DW는 알몸 엉금 분리). 등껍질은 들기·던지기. 초록: 절벽 낙하, 빨강: 방향 전환. 거대화 등껍질은 블록 파괴 & 초록/빨강. 알몸 엉금(Beach Koopa) & walk, spin(등껍질) \\
등껍질 Buzzy Beetle & 클래식 & 파이어볼 면역. 천장 부착 시 접근하면 낙하(회전 등껍질). 날개 시 파라비틀 & 버지 등껍질: 착용 가능(상부 방어) & walk, spin \\
가시 등껍질 Spike Top & 클래식 & 벽·천장 타고 순회. 파이어볼 면역 & 파랑 변형: 고속 & walk(경로) \\
가시돌이 Spiny & 전체 & 등에 가시, 밟기 불가. 천장 부착 시 버지처럼 행동. 너구리 꼬리·망토·고양이 할퀴기·부메랑으로 뒤집기 가능. 날개 시 가시 4방향 발사 & 가시 등껍질(착용, 적 처치 가능), 스파이니 에그 & walk \\
징오징오 Blooper & 전체 & 지상·수중 배치 가능, 플레이어 향해 불규칙 이동. 지상에서는 밟기 가능. 3DW에서는 벽 통과 불가 & 새끼 동반(Blooper Nanny) & swim \\
뽀꾸뽀꾸 Cheep Cheep & 전체 & 지상: 튀어다니며 도약 공격. 수중: 초록은 한 방향 직진, 빨강은 왕복+추적. 용암 배치 시 밟기 무효 & SMW: Blurp, NSMBU: Deep Cheep 외형 & swim / onair \\
스킵스퀵 Skipsqueak & 3DW & 제자리 달리기, 플레이어 점프에 맞춰 도약 & 가시 변형(밟기 불가) & onair \\
벌 Stingby & 3DW & 발견 시 수평으로만 추적 비행 & — & onair(비행) \\
뻐끔플라워 Piranha Plant & SMB1·SMB3·NSMBU·3DW & 토관 설치 시 주기 출입, 그 외 고정 배치. 3DW에서는 돌진하며 밟기 가능 & 파이어 뻐끔: 화염구 발사(전체). SMW: 점프 뻐끔 & idle, emerge \\
뻐끔 덩굴 Piranha Creeper & 3DW & 설정 경로 따라 신축. 보라: 상시 왕복, 파랑: 경로 끝에서 수면. 연타로 시작점까지 밀어내면 처치 & 보라/파랑 & 신축(경로) + idle \\
검은 뻐끔 Muncher & 클래식 & 고정 해저드. 구두·요시·등껍질로 안전 보행. POW 또는 하이힐 스톰프로만 처치. 날개 시 주기 도약 & — & idle \\
쿵쿵 Thwomp & 전체 & x축 접근 시 낙하 후 복귀. 착지로 ?블록·P스위치 작동. 가로 이동 설정 가능. 3DW: 접촉 안전, 압사만 즉사. 거대화 시 블록 연속 파괴. 밤 지상: 부끄부끄처럼 추적 & — & idle, 낙하(코드) \\
두더지 Monty Mole & 클래식 & 땅에 숨어 있다 접근 시 튀어나와 추적 & — & emerge, walk \\
렌치 두더지 Rocky Wrench & 클래식 & 맨홀에서 서서히 나와 렌치 투척 & — & emerge, attack \\
해머 브로스 Hammer Bro & 전체 & 아치형 해머 투척, 간헐 도약 & 슬레지 브로스(거대화): 착지 지진으로 경직 & idle, attack, onair \\
파이어 브로 Fire Bro [3.0] & 3DW & 화염구 투척 & 헤비(거대화): 지진 추가 & idle, attack \\
멍멍이 Chain Chomp & 클래식 & 말뚝에 묶여 돌진 반복. NSMBU 등에서 GP로 말뚝 제거 시 해방 & 사슬 없는 멍멍이(자유 이동) & idle(돌진 트윈) \\
가시 Spike [2.0] & 전체 & 가시공 생성·투척. SMB1·SMB3: 전방 비행+보행, 그 외: 굴리기+정지. 설원에서는 눈덩이 투척 & 가시공(Spike Ball)·눈덩이(단독 배치 가능, 밤엔 유도) & idle, attack \\
홉챱 Hop-Chops & 3DW & 추적. 밟으면 트램펄린으로 사용 가능 & — & onair \\
꿈틀이 Wiggler & 클래식 & 보행, 공격 안 함. 밟으면 분노(가속+추적). 날개 시 도약 보행. 밤 지상: 부유 & 분노 꿈틀이(처음부터 분노) & walk, 분노 walk \\
부끄부끄 Boo & 전체 & 등 돌리면 접근, 마주보면 정지. 지면 바로 위 배치 시 스트레치로 변함. 밤 지상: 행동 반전 & 부끄 무리(원형 회전, 틈 통과 가능) / 스트레치 & 2상태 idle \\
피파 Peepa & 3DW & 원형 궤도로 상시 이동 & 부끄부끄의 변형 & 궤도 이동 \\
용암방울 Lava Bubble & 전체 & 화면 하단에서 주기 도약. 날개 시 대각 비행 & — & onair \\
폭탄병 Bob-omb & 전체 & 보행. 공격받거나 불 접촉 시 점화, 들기·던지기 가능, 폭발로 블록 파괴. 날개 시 부유 & 점화 폭탄병(등장 즉시 점화, 약 5초 후 폭발) & walk, 점멸+폭발(코드) \\
드라이본 Dry Bones & 전체 & 밟으면 붕괴 후 재조립. SMW에서 뼈 투척. 파이어볼 면역. 거대화 시 스핀점프·GP 필요 & 드라이본 등껍질: 착용 시 용암 표류·죽은 척 무적 & walk, 붕괴/재조립 \\
피시본 Fish Bone & SMB1·SMB3·SMW·3DW & 수중에서 발견 시 직선 돌진, 벽 충돌 시 자괴 & — & swim, 돌진 \\
마귀 Magikoopa & 전체 & 순간이동하며 마법 발사. 마법이 블록을 다른 파츠로 변환 & 영국판 명칭 Kamek & idle, attack, teleport \\
포키 Pokey [2.0] & 전체 & 다단 몸통 보행, 접촉 피해. 높이 조절. 파이어볼·요시로 처치, 머리 타격 시 즉사. 밤 지상: 부유 & 스노 포키(설원): 밟으면 눈덩이로 축소 & walk(마디) \\
쿠파 Bowser & 클래식 & 도약+화염 브레스. SMB3: GP 경직, SMW·NSMBU: 공중 다연발 화염. 클라운 카 탑승 시 폭탄병 투척 & 3DW: 야옹쿠파(벽 오르기, 날개 시 배경에서 공격) & idle, attack, onair \\
쿠파주니어 Bowser Jr. & 클래식 & 화염 + 등껍질 돌진. 밟기 3회로 처치. 날개 시 해머 다연발 & — & idle, spin, attack \\
부메부메 Boom Boom & 전체 & 추적+도약, 공격받으면 가시 웅크림(클래식) / 회전 후 어지러움(3DW). 밟기 3회 & 3DW: 폼폼(분신+수리검, 본체는 분홍 수리검) & walk, spin, 기절 \\
메카쿠파 Mechakoopa [3.0] & 클래식 & 초록 엉금형 보행. 밟으면 기절, 들기·던지기. 방치 시 기상. 밤 지상: 제트 부유 & 블래스타(유도 미사일)·자파(수평 레이저) 변형 & walk, attack, 기절 \\
성난 태양 Angry Sun [2.0] & 클래식 & 화면을 돌다 급강하 공격. 구역당 1개 & 달(Moon): 밤 전환 + 접촉 시 화면 내 적 전멸 & idle, 급강하(코드) \\
김수한무 Lakitu & 클래식 & 상공 왕복, 가시돌이(변경 가능) 투하. 처치 시 구름 탑승 가능 & 라키투 구름(단독 배치, 시간제 탑승) & onair, attack \\
클라운 카 Koopa Clown Car & 클래식 & 탑승 비행체(무기한). 적을 태우면 플레이어 추적 & 파이어 클라운 카: 화염 발사(차지 가능). NSMBU: 주니어 클라운 카 & hover(코드) \\
엉금엉금 카 Koopa Troopa Car & 3DW & 알몸 엉금이 운전. 처치 후 탑승 운전. 3회 피격으로 파괴 & — & drive \\
가시복어 Porcupuffer & 3DW & 지상 테마: 수면 이동+도약해 삼키기 시도. 수중: 추적. 파이어볼 4발로 처치 & — & swim \\
차바아 Charvaargh & 3DW & 용암·측면·전경에서 아치형 도약 & — & 도약(코드 트윈) \\
불리 Bully & 3DW & 피해는 없으나 밀쳐서 경직·낙사 유도. 밀어서 절벽 밖으로 처치. 날개 시 낙하 직전 부유 & — & walk, 밀치기 \\
\midrule
\multicolumn{5}{l}{\textbf{쿠파 7인방 Koopalings [3.0] --- 클래식, 각각 밟기 3회. 날개 부착 시 개별 행동 변화}} \\
\midrule
래리 Larry & 클래식 & 도약하며 마법탄 발사 & 날개: 높은 점프 & idle, attack \\
이기 Iggy & 클래식 & 보행하며 고속 마법탄 & 날개: 도약 이동 & walk, attack \\
웬디 Wendy & 클래식 & 도약 + 벽·바닥 반사 링 발사 & 날개: 높은 점프 & idle, attack \\
레미 Lemmy & 클래식 & 공 위에서 균형 잡으며 바운스 볼 발사 & 날개: 공과 함께 도약 & idle, attack \\
로이 Roy & 클래식 & 마법탄 + 땅·천장 파고들어 출몰, 착지 지진 & 날개: 공중 체공 후 착지 & emerge, attack \\
모튼 Morton & 클래식 & 마법탄 + 착지 지진·지면 화염(벽돌 파괴) & 날개: 착지 페인트 & onair, attack \\
루드윅 Ludwig & 클래식 & 마법탄 3연발(피격 후 5연발) + 플러터 도약 & 날개: 체공 중 발사 & onair, attack \\
\bottomrule
\end{longtable}

% ============================================================
\section{장치 (Gizmos)}
% ============================================================

열쇠 아이콘, 노란색 카테고리. 총 40종(별도)·51종(고유 명칭).
대부분 코드 연출(회전·왕복·토글)로 표현 가능하며, AI 생성 애니메이션이
필요한 항목이 사실상 없다.

\begin{longtable}{L{2.9cm} L{1.4cm} L{6.2cm} L{2.7cm} L{1.8cm}}
\toprule
\textbf{파츠} & \textbf{스타일} & \textbf{특성} & \textbf{변형·비고} & \textbf{애니메이션} \\
\midrule
\endhead
버너 Burner & 클래식 & 주기적으로 화염 분사, 발판으로도 사용 가능. 방향 변경 가능 & 점선 표시 변형: 분사 타이밍 오프셋 & jet 루프(코드) \\
킬러 대포 Bill Blaster & 전체 & 킬러 무한 발사(플레이어 인접 시 정지). 높이 조절, 발사체 교체 가능. 3DW: 킬러가 벽 충돌 시 폭발(블록 파괴) & 유도탄 대포(빨강, Bull's-Eye): 유도 킬러 발사, 타 파츠는 더 멀리 발사 & 발사 플래시(코드) \\
매그넘 킬러 Banzai Bill & 전체 & 4방향 직선 비행 대형 킬러. 3DW: 배경에서 발사 가능, 블록 파괴, 고양이 할퀴기로 방향 전환 & 유도 매그넘(빨강). 3DW: 캣 매그넘 외형 & onair(직선) \\
대포 Cannon & 클래식 & 주기적으로 포탄 발사. 벽·천장 설치 가능 & 빨간 대포: 포탄 약 3배속 & 발사(코드) \\
고드름 Icicle & 전체 & 접근 시 떨림 후 낙하, 재생성. ?블록·P스위치 작동시킴. 위에서 밟기 가능, 아래 접촉 피해 & 고정 변형(낙하 없음, 진한 색) & shake + 낙하(코드) \\
회오리 Twister & 전체 & 플레이어·파츠를 상승 기류로 부양 & — & 소용돌이 루프 \\
!블록 ! Block & 3DW & 치면 설정 경로로 블록이 연장, 일정 시간 후 복귀. GP 시 즉시 전체 연장 & — & 연장(코드) \\
나무 Tree & 3DW & 덩굴처럼 오르기 가능. 아이템 수납(꼭대기 도달 시 방출). 높이 조절 & — & static \\
나무상자 Crate & 3DW & GP로 파괴, 들고 운반, 물·용암에 뜸. 코인 1개 수납 가능 & — & static \\
열쇠 Key & 전체 & 열쇠문·잠긴 워프 상자 개방. 핑크 코인 완수 시에도 출현 & 저주받은 열쇠 [3.0](SMB1): 획득 시 판토가 추격(동시 1마리) & idle 반짝임 \\
문 Warp Door & 전체 & 같은 구역 내 워프(쌍 배치) & P문: P스위치 중에만 출현. 열쇠문: 열쇠 필요 & 열림(코드) \\
워프 상자 Warp Box & 3DW & 접촉 시 워프, 진입 후 소멸 & 열쇠 워프 상자 & 반짝임(코드) \\
P블록 P Block [2.0] & 전체 & 파괴 불가 블록. P스위치 상태에 따라 실체$\leftrightarrow$비실체 전환 & 초기 상태 2종 & toggle \\
P스위치 P Switch & 전체 & 일시적으로 코인$\leftrightarrow$벽돌 전환, P블록 상태 전환, P문 출현. 천장 부착 가능 & — & 눌림(코드) \\
POW 블록 POW Block & 전체 & 타격 시 지상 적 전멸. 들기·던지기 가능 & 빨간 POW(3DW): 주변 벽돌 파괴 + 인접 블록 연쇄 작동 & shake(코드) \\
트램펄린 Trampoline & 전체 & 고반발 점프대. 운반 가능, 자체 중력 있음 & 가로 변형: 옆으로 튕김, 운반 불가 & 압축(코드) \\
덩굴 Vine & 클래식 & 오르기. 블록 수납 시 타격으로 성장 & — & 성장(코드) \\
화살표 표지판 Arrow Sign & 클래식 & 8방향 안내 표지 & — & static \\
중간 깃발 Checkpoint Flag & 전체 & 리스폰 지점. 구역당 1개. 파워업 수납 가능(통과 시 부여) & 클리어 조건과 병용 불가 & 깃발 상승(코드) \\
리프트 Lift & 클래식 & 수평/수직 왕복 발판. 길이 조절 & 낙하 리프트(Flimsy): 밟으면 추락, 레일 위에선 밟기 전 정지 & 이동(코드) \\
구름 리프트 Cloud Lift & 3DW & 다방향 이동 발판 & — & 이동(코드) \\
용암 리프트 Lava Lift & 클래식 & 밟으면 우측 이동, 시간 경과 후 침몰 & 고속 변형(파랑): 3배속, 침몰 없음 & float(코드) \\
시소 Seesaw & 클래식 & 하중 방향으로 기울어지는 발판. 길이 조절 & — & 물리(코드) \\
그라인더 Grinder & 클래식 & 회전 톱날 해저드. 다른 파츠는 통과 & — & rotate 루프(코드) \\
범퍼 Bumper & 클래식 & 접촉 시 튕겨냄. 위에서 밟으면 트램펄린처럼 사용 & — & 충돌 플래시(코드) \\
스큐어 Skewer & 클래식 & 갑자기 돌출하는 가시 기둥. 진로의 블록 첫 줄 파괴. 가로/세로 배치 & 영국판 명칭 Spike Pillar & 돌출 트윈(코드) \\
매달린 클로 Swinging Claw & 클래식 & 진자 운동, 잡고 스윙. 적·아이템을 매달아 접근 시 투하 가능 & — & swing(코드) \\
ON/OFF 스위치 & 전체 & 전역 상태 토글: 점선 블록·레일 분기·컨베이어 방향·가시 블록·ON/OFF 트램펄린 연동 & — & toggle \\
점선 블록 Dotted-Line Block & 전체 & 빨강: ON일 때 실체, 파랑: OFF일 때 실체 & 빨강/파랑 & toggle \\
가시 블록 Spike Block & 3DW & 가시가 돋았다 들어가는 대형 블록 & 노랑(교대) / 빨강(ON 시 가시) / 파랑(OFF 시 가시). 3.0에서 지형$\to$장치로 이동 & toggle \\
스네이크 블록 Snake Block & 전체 & 밟으면 설정 경로 따라 이동하는 블록 열 & 고속 변형(파랑, 약 2배속) & 경로 이동(코드) \\
파이어바 Fire Bar & 클래식 & 블록 축 중심 회전 화염 막대. 시작 각도·길이·방향 설정 & — & rotate 루프(코드) \\
일방통행벽 One-Way Wall & 클래식 & 한 방향으로만 통과 & — & static \\
컨베이어 벨트 Conveyor Belt & 전체 & 위 개체를 설정 방향으로 이동. 경사 배치 가능, ON/OFF로 방향 전환 & 고속 변형(약 2배속) & 벨트 스크롤(코드) \\
레일 Track & 클래식 & 파츠 이동 경로. 대각 레일 흔들면 곡선. ON/OFF로 분기 전환 & — & static \\
레일 블록 Track Block & 3DW & 경로 따라 이동하는 블록. 빨강: 상시, 파랑: 밟으면 시작(종점 도달 후 복귀) & 빨강/파랑 & 경로 이동(코드) \\
ON/OFF 트램펄린 [3.0] & 3DW & 스위치 상태에 따라 트램펄린$\leftrightarrow$일반 발판 & — & toggle \\
버섯 트램펄린 Mushroom Trampoline & 3DW & 고반발 버섯 기둥. 주황: 왕복 이동, 파랑: 고정 & 주황/파랑 & 압축(코드) \\
대시 블록 Dash Block [2.0] & 3DW & 밟으면 대시 패널처럼 가속 부여 & — & static \\
점멸 블록 Blinking Block & 3DW & 분홍/파랑이 리듬에 맞춰 교대로 실체화 & — & toggle(박자) \\
\bottomrule
\end{longtable}

% ============================================================
\section{전역 요소 (파츠 외)}
% ============================================================

\subsection{게임 스타일과 테마}

게임 스타일 5종: SMB1 / SMB3 / SMW / NSMBU (상호 전환 가능) + 3D월드
(전환 시 에디터 초기화되는 별도 계열). 테마는 10종이며 각각 낮/밤 버전이
있다(밤은 성난 태양$\to$달 전환으로 최초 해금).

\begin{longtable}{L{2.2cm} L{13.2cm}}
\toprule
\textbf{테마} & \textbf{밤 효과 (원문 근거)} \\
\midrule
지상 & 적·아이템의 행동이 다양하게 변화 (1UP$\to$썩은 버섯, 굼바 부유 등) \\
지하 & 화면 상하 반전 + 수직 조작 반전 \\
수중 & 플레이어·일부 요소 주변 외 화면 암전 \\
유령의 집 & \textcolor{red}{원문 표에 효과 기재 없음 (확인 필요)} \\
비행선 & 배경 폭풍(SMB1 제외) + 전 요소가 수중처럼 거동(가시돌이는 감속) \\
성 & 플레이어만 수중처럼 거동, 다른 파츠는 불변 \\
사막 & 스타일별로 다른 모래폭풍 발생 \\
설원 & 모든 표면 미끄러움 \\
숲 & 물이 독물로 변화 \\
하늘 & 저중력 + 전 요소 수중 거동(가시돌이는 감속) \\
\bottomrule
\end{longtable}

\subsection{에디터 설정}

\begin{itemize}
  \item \textbf{제한시간}: 10--500초, 10초 단위
  \item \textbf{자동 스크롤(Custom Scroll)}: 3단 속도, 구역당 방향 전환 최대 10회
  \item \textbf{스크롤 스톱}: 상하 전체를 블록으로 막으면 카메라 정지
  \item \textbf{수위}: 숲(물)·성(용암)·밤 숲(독물)에서 시작 높이와 상승/승강 주기 설정 (하강만은 불가)
  \item \textbf{서브 에어리어}: 가로 또는 세로(최대 164블록) 보조 맵
  \item \textbf{클리어 조건}: 레벨당 최대 1개. 세 분류 --- 행동(무피해 도달, 무착지 도달) / 파츠(특정 적 $n$마리 처치 후 도달, 코인류 $n$개 수집 후 도달, POW·P스위치 $n$개 작동 후 도달 등) / 상태(특정 폼·탑승·소지 상태로 도달). 중간 깃발과 병용 불가
\end{itemize}

\subsection{파츠 개수 제한 (구역당)}

\begin{longtable}{L{6cm} L{2.5cm}}
\toprule
\textbf{대상} & \textbf{제한} \\
\midrule
땅 & 4000 \\
블록류 & 2000 \\
레일 & 1500 \\
고드름 & 300 \\
발판류 & 200 \\
화살표/일방통행벽 & 100 \\
적·아이템 / 파워업 / 투명 토관 & 각 100 \\
엉금 카 / 뻐끔 덩굴 / 레일 블록 / !블록 & 각 10 \\
스네이크 블록 / 부메부메 / 핑크 코인 / 대형 코인 & 각 5 \\
문 / 워프 상자 & 각 4 \\
쿠파 / 차바아 & 각 3 \\
성난 태양(달) / 중간 깃발 & 각 1 \\
사운드 이펙트 & 300 \\
\bottomrule
\end{longtable}

\subsection{스토리 모드 전용 요소}

돌(Stone, 운반용 무거운 블록)과 토드(Toad, 골까지 호위) 2종은 코스 메이커에서
사용 불가. 사운드 이펙트(감정/스팅어/리액션/연출/음악 5분류)는 게임플레이에
영향이 없어 본 카탈로그에서 세부 생략.

% ============================================================
\section{우리 프로젝트로의 매핑 (팀 설계안)}
% ============================================================

\subsection{카테고리 매핑}

\begin{longtable}{L{2.8cm} L{3.4cm} L{9.2cm}}
\toprule
\textbf{우리 카테고리} & \textbf{SMM2 대응} & \textbf{비고} \\
\midrule
플랫폼 & 지형 & 유저 제작 대상. 속성 조합으로 블록류 재현 \\
몬스터 & 적 & 유저 제작 대상. idle/walk/onair 적용 \\
장애물 & 지형 일부 + 장치 & 유저 제작 대상. 함정류를 속성으로 재현 \\
아이템 & 아이템 & (도입 시) 코인·1회용 파워업 정도로 축소 권장 \\
고정 요소 & 시작/골/중간 깃발 & 시스템 제공. 유저는 위치만 지정 \\
\bottomrule
\end{longtable}

\subsection{속성 슬롯 초안 (검증된 원작 파츠에서 역산)}

\paragraph{플랫폼 속성}
\begin{itemize}
  \item 충돌 방식: 완전 충돌 / 반통과 --- (땅 vs 반통과 발판·다리)
  \item 밟으면 $n$초 후 낙하, $m$초 후 재생 (T/F, $n$, $m$) --- (도넛 블록, 고드름의 재생성)
  \item 파괴 가능 (T/F) --- (벽돌 블록, 얼음 블록)
  \item 표면 마찰: 보통 / 미끄러움 --- (얼음 블록, 설원 밤 효과)
  \item 밟으면 튕김 (반발력 $n$) --- (음표 블록, 트램펄린, 버섯 트램펄린)
  \item 왕복/경로 이동 (경로, 속도) --- (리프트, 스네이크 블록, 레일 블록)
\end{itemize}

\paragraph{몬스터 속성}
\begin{itemize}
  \item 생명력 $n$ --- (부메부메·쿠파주니어·7인방의 3회 밟기, 가시복어의 파이어볼 4발)
  \item 감지 시 활성화 (T/F, 거리 $n$) --- (3DW 굼바의 발견 시 추적, 두더지, 쿵쿵)
  \item 이동 유형: 고정 / 보행 / 비행 / 추적 --- (검은 뻐끔 / 굼바 / 킬러 / 부끄부끄·벌)
  \item 절벽에서 방향 전환 (T/F) --- (굼바 vs Goombrat, 초록 vs 빨강 엉금)
  \item 밟기 가능 (T/F) --- (굼바 vs 가시돌이)
  \item 이동 속도 $n$ --- (가시 등껍질 빨강 vs 파랑, 고속 변형 전반)
\end{itemize}

\paragraph{장애물 속성}
\begin{itemize}
  \item 닿으면 죽음/피해 (T/F) --- (가시, 그라인더, 파이어바)
  \item 플레이어와 $x$(또는 $y$)좌표 겹치면 낙하 (T/F) --- (쿵쿵, 고드름)
  \item 주기 $n$초 돌출/분사 (T/F, $n$) --- (스큐어, 버너, 가시 블록)
  \item 회전 (T/F, 속도) --- (파이어바, 그라인더)
\end{itemize}

\subsection{애니메이션 체계 결론}

검증된 전 파츠를 통틀어 개체 고유 애니메이션은 idle / walk(swim 포함) /
onair 계열 루프가 전부이며, 나머지(bump·shake·rotate·toggle·경로 이동·
돌출·낙하)는 모두 코드 연출로 원작이 처리하는 영역이다. 따라서 확정한
3종 체계(idle/walk/onair)로 원작급 표현이 가능하다. 몬스터의 이동 유형이
비행·수중인 경우 walk 대신 onair/swim 프롬프트를 재사용한다.

\end{document}
